%% maths-nombre-derive — construction du nombre dérivé f'(a) comme limite du taux
%% d'accroissement. Courbe C_f de f(x) = 0,20 x^2 + 0,40 ; A d'abscisse a = 1,2 et
%% M d'abscisse a + h = 3,4.
%%   - sécante (AM) en trait fin bleu : sa pente est le taux d'accroissement ;
%%   - tangente en A en trait épais rouge : sa pente est le nombre dérivé ;
%%   - le triangle rectangle A-P-M porte les deux cotes h et f(a+h) - f(a) ;
%%   - la flèche en tirets rappelle que M glisse vers A quand h tend vers 0.
%% Les DEUX formules (taux d'accroissement et nombre dérivé) sont dans l'encadré sous
%% la figure, enveloppé de \rappel : la variante « maths-nombre-derive-muet » ne montre
%% que la construction géométrique (à utiliser dans les énoncés).
\documentclass[tikz,border=4pt]{standalone}
\usepackage{liaisons}
\begin{document}
\begin{tikzpicture}[x=1.25cm,y=1.05cm,
  axe/.style={-{Stealth[length=5pt]}, line width=0.6pt},
  courbe/.style={line width=1.3pt, line join=round, black!80},
  secante/.style={solideB, line width=0.7pt},
  tangente/.style={solideA, line width=1.6pt},
  cote/.style={solideE, line width=0.9pt},
  pointille/.style={line width=0.5pt, dash pattern=on 3pt off 2pt, black!50},
  etiq/.style={font=\small, inner sep=2pt}]

%% ---------------- la fonction et les deux abscisses -----------------------
\def\aa{1.2}      % abscisse de A
\def\hh{2.2}      % accroissement h
\def\fa{0.688}    % f(a)
\def\fah{2.712}   % f(a+h)
\def\ps{0.92}     % pente de la sécante (AM)
\def\pt{0.48}     % pente de la tangente en A

%% ---------------- axes -----------------------------------------------------
\draw[axe] (-0.45,0) -- (4.35,0) node[anchor=north east, inner sep=3pt, font=\small] {$x$};
\draw[axe] (0,-0.35) -- (0,3.95) node[anchor=south east, inner sep=3pt, font=\small] {$y$};
\node[etiq, anchor=north east] at (-0.03,-0.03) {$O$};

%% ---------------- la courbe ------------------------------------------------
\draw[courbe] plot[domain=-0.35:4.05, samples=90, smooth] (\x,{0.20*\x*\x+0.40});
\node[font=\small, text=black!80, anchor=west, inner sep=3pt] at (4.05,3.7) {$\mathcal{C}_f$};

%% ---------------- traits de rappel vers les axes ---------------------------
\draw[pointille] (\aa,0) -- (\aa,\fa) -- (0,\fa);
\draw[pointille] (\aa+\hh,0) -- (\aa+\hh,\fah) -- (0,\fah);
\node[etiq, anchor=north] at (\aa,-0.04) {$a$};
\node[etiq, anchor=north] at (\aa+\hh,-0.04) {$a+h$};
\node[etiq, anchor=east] at (-0.05,\fa) {$f(a)$};
\node[etiq, anchor=east] at (-0.05,\fah) {$f(a+h)$};

%% ---------------- sécante (AM) et tangente en A ----------------------------
\draw[secante] (0.35,{\fa+\ps*(0.35-\aa)}) -- (3.95,{\fa+\ps*(3.95-\aa)});
\node[font=\small, text=solideB, anchor=south east, inner sep=2pt, rotate=41]
     at (2.95,{\fa+\ps*(2.95-\aa)}) {sécante $(AM)$};
\draw[tangente] (-0.3,{\fa+\pt*(-0.3-\aa)}) -- (2.95,{\fa+\pt*(2.95-\aa)});
\node[font=\small, text=solideA, anchor=north east, inner sep=2pt, rotate=24]
     at (2.92,{\fa+\pt*(2.92-\aa)-0.02}) {tangente en $A$};

%% ---------------- triangle du taux d'accroissement -------------------------
\draw[cote] (\aa,\fa) -- (\aa+\hh,\fa);
\draw[cote] (\aa+\hh,\fa) -- (\aa+\hh,\fah);
\draw[line width=0.5pt, black!60] (\aa+\hh-0.16,\fa) -- (\aa+\hh-0.16,\fa+0.14) -- (\aa+\hh,\fa+0.14);
\node[font=\small, text=solideE, anchor=north, inner sep=3pt] at ({\aa+\hh/2},\fa) {$h$};
\node[font=\small, text=solideE, anchor=west, inner sep=3pt] at (\aa+\hh,{(\fa+\fah)/2})
     {$f(a+h)-f(a)$};

%% ---------------- points A et M --------------------------------------------
\fill[solideA] (\aa,\fa) circle (2pt);
\node[etiq, anchor=south east] at (\aa,\fa) {$A$};
\fill[solideB] (\aa+\hh,\fah) circle (2pt);
\node[etiq, anchor=south west, inner sep=3pt] at (\aa+\hh,\fah) {$M$};

%% ---------------- M glisse vers A (sous l'axe, pour ne rien recouvrir) -----
\draw[-{Stealth[length=6pt]}, line width=0.7pt, black!55, dash pattern=on 3pt off 2pt]
     (\aa+\hh,-0.42) to[bend left=22] (\aa+0.08,-0.42);
\node[font=\small, text=black!60, anchor=north, inner sep=2pt] at (2.3,-0.86)
     {quand $h\to0$, $M$ glisse vers $A$ le long de $\mathcal{C}_f$};

%% ---------------- ce que l'exercice demande de retrouver -------------------
\rappel{\node[draw=black!45, fill=black!3, rounded corners=3pt, line width=0.5pt,
      inner sep=5pt, anchor=north, align=center, font=\small] at (2.0,-1.42)
     {pente de la sécante $(AM)$ $=$ taux d'accroissement :\\[2pt]
      $\dfrac{f(a+h)-f(a)}{h}$\\[5pt]
      pente de la tangente en $A$ $=$ nombre dérivé :\\[2pt]
      $f'(a)=\displaystyle\lim_{h\to0}\dfrac{f(a+h)-f(a)}{h}$};}
\end{tikzpicture}
\end{document}
